\documentclass[a4paper,11pt]{book}
%\documentclass[a4paper,twoside,11pt,titlepage]{book}

\usepackage[
  a4paper,
  inner=35mm,      % margen interior (libro a doble cara)
  outer=30mm,      % margen exterior
  top=25mm,        % margen superior real
  bottom=30mm,     % margen inferior
  headheight=10pt, % altura de la cabecera para fancyhdr
  headsep=15pt,    % separación cabecera–texto (reduce “aire” arriba)
  includehead      % cuenta la cabecera dentro del margen superior
]{geometry}


\usepackage{listings}
\usepackage[utf8]{inputenc}
\usepackage[spanish]{babel}

\usepackage{csquotes}

\usepackage[
  style = apa,
  backend = biber,
  ]
{biblatex}
\addbibresource{bibliografia.bib}
%\bibliography{bibliografia/bibliografia}\addcontentsline{toc}{chapter}{Bibliografía}

% \usepackage[style=list, number=none]{glossary} %
%\usepackage{titlesec}
%\usepackage{pailatino}

\decimalpoint
\usepackage{dcolumn}
\newcolumntype{.}{D{.}{\esperiod}{-1}}
\makeatletter
\addto\shorthandsspanish{\let\esperiod\es@period@code}
\makeatother

%para los colores de las tablas
\usepackage[table,xcdraw]{xcolor}

%\usepackage[chapter]{algorithm}
\RequirePackage{verbatim}
%\RequirePackage[Glenn]{fncychap}
\usepackage{fancyhdr}
\usepackage{graphicx}
\usepackage{afterpage}
\usepackage{caption}
\usepackage{float}
\usepackage{pdflscape}
\usepackage{longtable}
\usepackage{amsmath}

\usepackage[pdfborder={000}]{hyperref} %referencia

\DeclareCiteCommand{\citep}[\mkbibparens]
  {\usebibmacro{prenote}}
  {\usebibmacro{citeindex}\usebibmacro{cite}}
  {\multicitedelim}
  {\usebibmacro{postnote}}

% ********************************************************************
% Re-usable information
% ********************************************************************
\newcommand{\director}{Daniel Bolaños Martínez}
\newcommand{\otrodirector}{Alberto Durán López}
\newcommand{\thetitle}{Título del Trabajo}
\newcommand{\thetitleen}{Title in english}
\newcommand{\theauthor}{Mario Rosas Navarro}
\newcommand{\curso}{2025-2026}
\newcommand{\dni}{77650905FF}
\newcommand{\grado}{Doble Grado en Ingeniería Informática - Administración y Dirección de empresas}
\newcommand{\myFaculty}{Escuela Técnica Superior de Ingenierías Informática y de
Telecomunicación}
\newcommand{\myFacultyShort}{E.T.S. de Ingenierías Informática y de
Telecomunicación}
\newcommand{\myUni}{\protect{Universidad de Granada}}
\newcommand{\myLocation}{Granada a \today.}
\newcommand{\myVersion}{Version 0.1}


\hypersetup{
pdfauthor = {\theauthor (email (en) ugr (punto) es)},
pdftitle = {\thetitle},
pdfsubject = {},
pdfkeywords = {palabra_clave1, palabra_clave2, palabra_clave3, ...},
pdfcreator = {LaTeX con el paquete ....},
pdfproducer = {pdflatex}
}

%\hyphenation{}


%\usepackage{doxygen/doxygen}
%\usepackage{pdfpages}
\usepackage{url}
\usepackage[type={CC}, modifier={by-sa}, version={4.0}]{doclicense}
\usepackage{hyperref}
\usepackage{hyperxmp}
\usepackage{colortbl,longtable}
\usepackage[stable]{footmisc}
%\usepackage{index}

%\makeindex
%\usepackage[style=long, cols=2,border=plain,toc=true,number=none]{glossary}
% \makeglossary

% Definición de comandos que me son tiles:
%\renewcommand{\indexname}{Índice alfabético}
%\renewcommand{\glossaryname}{Glosario}

\pagestyle{fancy}
\fancyhf{}
\fancyhead[LO]{\leftmark}
\fancyhead[RE]{\rightmark}
\fancyhead[RO,LE]{\textbf{\thepage}}
\renewcommand{\chaptermark}[1]{\markboth{\textbf{#1}}{}}
\renewcommand{\sectionmark}[1]{\markright{\textbf{\thesection. #1}}}

%\setlength{\headheight}{1.5\headheight}

\newcommand{\HRule}{\rule{\linewidth}{0.5mm}}
%Definimos los tipos teorema, ejemplo y definición podremos usar estos tipos
%simplemente poniendo \begin{teorema} \end{teorema} ...
\newtheorem{teorema}{Teorema}[chapter]
\newtheorem{ejemplo}{Ejemplo}[chapter]
\newtheorem{definicion}{Definición}[chapter]

\addto\captionsspanish{%
  \renewcommand{\tablename}{Tabla}%
  \renewcommand{\listtablename}{Índice de Tablas}%
  \renewcommand{\listfigurename}{Índice de Figuras}%
}

\definecolor{gray97}{gray}{.97}
\definecolor{gray75}{gray}{.75}
\definecolor{gray45}{gray}{.45}
\definecolor{gray30}{gray}{.94}

\lstset{ frame=Ltb,
     framerule=0.5pt,
     aboveskip=0.5cm,
     framextopmargin=3pt,
     framexbottommargin=3pt,
     framexleftmargin=0.1cm,
     framesep=0pt,
     rulesep=.4pt,
     backgroundcolor=\color{gray97},
     rulesepcolor=\color{black},
     %
     stringstyle=\ttfamily,
     showstringspaces = false,
     basicstyle=\scriptsize\ttfamily,
     commentstyle=\color{gray45},
     keywordstyle=\bfseries,
     %
     numbers=left,
     numbersep=6pt,
     numberstyle=\tiny,
     numberfirstline = false,
     breaklines=true,
   }
 
% minimizar fragmentado de listados
\lstnewenvironment{listing}[1][]
   {\lstset{#1}\pagebreak[0]}{\pagebreak[0]}

\lstdefinestyle{CodigoC}
   {
	basicstyle=\scriptsize,
	frame=single,
	language=C,
	numbers=left
   }
\lstdefinestyle{CodigoC++}
   {
	basicstyle=\small,
	frame=single,
	backgroundcolor=\color{gray30},
	language=C++,
	numbers=left
   }

 
\lstdefinestyle{Consola}
   {basicstyle=\scriptsize\bf\ttfamily,
    backgroundcolor=\color{gray30},
    frame=single,
    numbers=none
   }


\newcommand{\bigrule}{\titlerule[0.5mm]}


%Para conseguir que en las páginas en blanco no ponga cabecerass
\makeatletter
\def\clearpage{%
  \ifvmode
    \ifnum \@dbltopnum =\m@ne
      \ifdim \pagetotal <\topskip
        \hbox{}
      \fi
    \fi
  \fi
  \newpage
  \thispagestyle{empty}
  \write\m@ne{}
  \vbox{}
  \penalty -\@Mi
}
\makeatother
\usepackage{pdfpages}
\begin{document}
\begin{titlepage}
 
\newlength{\centeroffset}
\setlength{\centeroffset}{-0.5\oddsidemargin}
\addtolength{\centeroffset}{0.5\evensidemargin}
\thispagestyle{empty}

\noindent\hspace*{\centeroffset}\begin{minipage}{\textwidth}

\centering
\includegraphics[width=0.9\textwidth]{imagenes/logo_ugr.jpg}\\[1.4cm]

\textsc{ \Large TRABAJO FIN DE GRADO\\[0.2cm]}
\textsc{\large \grado}\\[0.2cm]
\textsc{Curso Académico \curso}\\[1cm]
% Upper part of the page
% 
% Title
{\huge\bfseries \thetitle}
\end{minipage}

\vspace{2.5cm}
\noindent\hspace*{\centeroffset}\begin{minipage}{\textwidth}
\centering

\textbf{Autor}\\ {\theauthor}\\[2.5ex]
\textbf{Directores}\\
{\director \\ \otrodirector}\\[2cm]

\vspace{2cm}

\includegraphics[width=0.3\textwidth]{imagenes/etsiit_logo.png}\\[0.1cm]
\textsc{Escuela Técnica Superior de Ingenierías Informática y de Telecomunicación}\\
\textsc{---}\\
Granada, a \today
\end{minipage}

\end{titlepage}



\begin{titlepage}
 
 
\setlength{\centeroffset}{-0.5\oddsidemargin}
\addtolength{\centeroffset}{0.5\evensidemargin}
\thispagestyle{empty}

\noindent\hspace*{\centeroffset}\begin{minipage}{\textwidth}

\centering

 \vspace{3.3cm}

%si el proyecto tiene logo poner aquí
\includegraphics[width=0.25\linewidth]{imagenes/logo proyecto.png} 
 \vspace{0.5cm}

% Title
{\huge\bfseries \thetitle}
\end{minipage}

\vspace{2.5cm}
\noindent\hspace*{\centeroffset}\begin{minipage}{\textwidth}
\centering

\textbf{Autor}\\ {\theauthor}\\[2.5ex]
\textbf{Directores}\\
{\director \\ \otrodirector}\\[2cm]
%\includegraphics[width=0.15\textwidth]{imagenes/tstc.png}\\[0.1cm]
%\textsc{Departamento de Teoría de la Señal, Telemática y Comunicaciones}\\
%\textsc{---}\\
%Granada, mes de 201
\end{minipage}
%\addtolength{\textwidth}{\centeroffset}
\vspace{\stretch{2}}

 
\end{titlepage}



\chapter*{}

\thispagestyle{empty}

\begin{center}
{\large\bfseries \thetitle}\\
\end{center}
\begin{center}
\theauthor\\
\end{center}

\noindent{\textbf{Resumen}}\\

\vspace{0.5cm}

Las smart cities utilizan tecnologías de la información y comunicación para mejorar la calidad de vida de sus habitantes. Un componente esencial de estas ciudades inteligentes es la capacidad de capturar y procesar grandes volúmenes de datos provenientes de diversas fuentes, como sensores, dispositivos IoT (Internet of Things), sistemas de transporte, servicios públicos y redes sociales. Estos datos se utilizan para monitorear y gestionar los recursos urbanos de manera eficiente.

Para llevar a cabo esta monitorización, se emplean dashboards. Los cuadros de mando, paneles de control o dashboards son herramientas fundamentales para la gestión y supervisión de distintos aspectos de las smart cities. Estos paneles ofrecen visualizaciones de datos en tiempo real, lo que facilita a los responsables de la toma de decisiones la gestión de recursos, la optimización de servicios y la mejora de la calidad de vida de los ciudadanos.

En este trabajo se presenta el análisis, diseño e implementación de una solución integral para la monitorización de estos datos, cumpliendo con los requisitos de utilizar software libre, asegurar que el sistema esté contenerizado para garantizar su desacoplamiento, y emplear un orquestador de contenedores que asegure una alta disponibilidad.

\vspace{0.7cm}
\noindent{\textbf{Palabras clave}: Cuadro de mandos, Base de datos, ETL, Orquestación de contenedores, Kubernetes, Software Libre, Monitorización}\\

\cleardoublepage
\chapter*{}

\thispagestyle{empty}

\begin{center}
{\large\bfseries \thetitleen}\\
\end{center}
\begin{center}
\theauthor\\
\end{center}

\vspace{0.5cm}
\noindent{\textbf{Abstract}}\\

Smart cities use information and communication technologies to improve the quality of life of their inhabitants. An essential component of these intelligent cities is the ability to capture and process large volumes of data from various sources, such as sensors, Internet of Things (IoT) devices, transportation systems, public services, and social networks. This data is used to monitor and manage urban resources efficiently.

To carry out this monitoring, dashboards are employed. Dashboards, also known as control panels or command boards, are essential tools for the management and supervision of different aspects of smart cities. These panels provide real-time data visualizations, which help decision-makers manage resources, optimize services, and improve the quality of life for citizens.

This work presents the analysis, design, and implementation of a comprehensive solution for monitoring this data, meeting the requirements of using open-source software, ensuring that the system is containerized to guarantee decoupling, and employing a container orchestrator to ensure high availability.

\vspace{0.7cm}
\noindent{\textbf{Keywords}: Dashboard, Database, ETL, Container Orchestration, Kubernetes, Open Source Software, Monitoring}\\

\chapter*{}
\thispagestyle{empty}

\noindent\rule[-1ex]{\textwidth}{2pt}\\[4.5ex]

D. \director, Profesor del Departamento Lenguajes y Sistemas de la Universidad de Granada. \\

D. \otrodirector, Profesor del Departamento Lenguajes y Sistemas de la Universidad de Granada. \\

\vspace{0.5cm}

\textbf{Informa:} \\

Que el presente trabajo, titulado \thetitle, ha sido realizado bajo su supervisión por \theauthor, y autorizamos la defensa de dicho trabajo ante el tribunal que corresponda.

\vspace{1cm}

Y para que conste, expiden el presente informe en \myLocation

\chapter*{}
\thispagestyle{empty}

\noindent\rule[-1ex]{\textwidth}{2pt}\\[4.5ex]

\textbf{Declaración de Originalidad} \\

Yo, \textbf{\theauthor}, alumno de la titulación \grado \ de la \myFacultyShort \ de la \myUni, con DNI \dni, declaro explícitamente que el trabajo presentado, correspondiente al curso académico \curso, es original, entendida esta, en el sentido de que no ha utilizado para la elaboración del trabajo fuentes sin citarlas debidamente.

\vspace{4cm}

\noindent Fdo: \theauthor

\vspace{1cm}

\begin{flushright}
\myLocation
\end{flushright}


\chapter*{}
\thispagestyle{empty}

\noindent\rule[-1ex]{\textwidth}{2pt}\\[4.5ex]

Yo, \textbf{\theauthor}, alumno de la titulación \grado \ de la \myFacultyShort \ de la \myUni, con DNI \dni, autorizo la ubicación de la siguiente copia de mi Trabajo Fin de Grado en la biblioteca del centro para que pueda ser consultada por las personas que lo deseen.

\vspace{4cm}

\noindent Fdo: \theauthor

\vspace{1cm}

\begin{flushright}
\myLocation
\end{flushright}

\vspace{3cm}

\doclicenseThis
The source code of this text and developed programs are available in the Github repository \url{https://github.com/}


\chapter*{Agradecimientos}
\thispagestyle{empty}

       \vspace{1cm}


Agradecimientos escritos...


\frontmatter
\tableofcontents
\listoffigures
\listoftables

\mainmatter
\setlength{\parskip}{5pt}

\input{capitulos/01_Introduccion}
\input{capitulos/02_Marco_teorico}
\input{capitulos/03_Estado_del_arte}
\input{capitulos/04_Tecnologias}
\input{capitulos/05_Diseño_Desarrollo}
\input{capitulos/06_Desarrollo_Futuro}

%%% --- The following two lines are what needs to be added --- %%%
\setcounter{biburllcpenalty}{7000}
\setcounter{biburlucpenalty}{8000}
\nocite{*}
\printbibliography[
title={Bibliografía},
heading=bibintoc]
%\bibliography{bibliografia/bibliografia}\addcontentsline{toc}{chapter}{Bibliografía}
%\bibliographystyle{unsrt}


%
%\appendix
%\input{apendices/manual_usuario/manual_usuario}
%%\input{apendices/paper/paper}
%\input{glosario/entradas_glosario}
% \addcontentsline{toc}{chapter}{Glosario}
% \printglossary

\end{document}
